\documentclass{article}
\usepackage[utf8]{inputenc}

\title{Reply to Editor and Reviewers}
\author{}
\date{\today}


%\usepackage{mwe}
%\usepackage{blindtext}
\usepackage{booktabs}
\usepackage{tabularx}  % in preamble

\usepackage{enumitem}
 
\usepackage{graphicx}
\usepackage{caption} % to get nicer top captions with tables
\usepackage{subcaption}

\newcommand{\mAP}{\textrm{mAP}}
\newcommand{\mysubsection}[1]{\subparagraph{#1}}

% to change the colour of text
\usepackage{xcolor}
% to strike through text
\usepackage[normalem]{ulem}
% to frame text
\usepackage{framed}

% to help with quoting added material in the revised paper 
\usepackage{quoting}
% \newcommand{\say}[1]{%
% \begin{quoting}
% #1
% \end{quoting}
% }

\newcommand{\sayold}[1]{%
{\color{olive}#1}
}



\newcommand{\saydel}[1]{%
{\color{darkgray}\sout{#1}}}

\newcommand{\saymod}[1]{%
  \texorpdfstring{\textcolor{blue}{#1}}{#1}%
}

% \newcommand{\saymod}[1]{#1}


\newcommand{\sayintro}[1]{%
  \texorpdfstring{\textcolor{olive}{#1}}{#1}%
}

\newcommand{\saynew}[1]{%
{\color{orange}#1}}

\newcommand{\reviewcomment}[3]{%
\vspace{4mm}
\textbf{C[$R_{#1}C_{#2}$]:} {\it #3}
\vspace{2mm}
}

\newcommand{\response}[1]{%
{\textbf{R:}} {#1}
}

% Change the default figure/table numbering to avoid confusion with main manuscript
\renewcommand{\figurename}{Fig.}
\renewcommand{\thefigure}{L\arabic{figure}}
\renewcommand{\thetable}{L\arabic{table}} 
\newcommand{\figref}[1]{\figurename~\ref{#1}}
\newcommand{\tabref}[1]{\tablename~\ref{#1}}
\newcommand{\subresponse}[1]{\underline{\MakeUppercase{#1}}}
% \newcommand{\enclose}[1]{\par\indent """  #1   \indent """\par}

\newcommand{\enclose}[1]{%
%    \smallskip
%    \par
%    \noindent%
%    \fbox{%
%        \begin{minipage}{\dimexpr\linewidth-2\fboxsep-2\fboxrule\relax}#1 
%        \end{minipage}%
%    }%
%    \par
\begin{framed}
#1
\end{framed}
}

\usepackage[pdfusetitle,hidelinks]{hyperref}

\begin{document}

\maketitle
%%%
\section*{Response to Editor}

We thank the Editor and reviewers for their positive assessment of our manuscript and for the helpful suggestions. We have revised the manuscript to address all remaining editorial comments. All updated references are marked with a magenta note, \textcolor{magenta}{updated}, in the revised manuscript.

\vspace{4mm}

\textbf{Editor Comment 1:} 
Please replace any reference to an old arXiv, which has not been reviewed, verified and published according to an accepted scientific review process for journal papers, with an appropriate reference from a peer reviewed journal. This applies in particular to References \#17, 18, 19 and 21.

\textbf{Response:} 
We thank the editor for this suggestion and have carefully revised the references accordingly. Where peer-reviewed versions were available, we have replaced the original arXiv citations with their corresponding published versions. The following reference has been updated to a peer-reviewed publication.

\begin{itemize}
\item[21] Loshchilov, I., Hutter, F.: Decoupled weight decay regularization. In: International Conference on Learning Representations (ICLR) (2019). 
\end{itemize}

However, for several references, no peer-reviewed publication is currently available. These works correspond to widely used datasets, benchmarking protocols, or software frameworks that are commonly cited in the literature. In these cases, the original arXiv references were retained to ensure proper attribution and reproducibility:

\begin{itemize}
\item[17] Cholecseg8k: a semantic segmentation dataset for laparoscopic cholecystectomy based on cholec80. (no peer-reviewed version currently available)
\item[18] Data Splits and Metrics for Method Benchmarking on Surgical Action Triplet Dataset (no peer-reviewed version currently available)
\item[19] MMDetection: Open mmlab detection toolbox and benchmark. (no peer-reviewed version currently available)
\end{itemize}

\vspace{4mm}
\textbf{Editor Comment 2:} 
Please complete any reference that contains an et al with the appropriate lists of authors. This applies in particular to References \#3, 11 and 13.

\textbf{Response:} 
We have updated the affected references to include the full list of authors in accordance with the journal guidelines.

\begin{itemize}
\item[3] Yamlahi, A., Tran, T.N., Godau, P., Schellenberg, M., Michael, D., Smidt, F.-H., N¨olke, J.-H., Adler, T.J., Tizabi, M.D., Nwoye, C.I., Padoy, N., Maier-Hein, L.: Self-distillation for surgical action recognition. In: Medical Image Computing and Computer Assisted Intervention – MICCAI 2023, pp. 637–646. Springer, Cham (2023).
\item[11] Younis, R., Yamlahi, A., Bodenstedt, S., Scheikl, P.M., Kisilenko, A., Daum, M., Schulze, A., Wise, P.A., Nickel, F., Mathis-Ullrich, F., Maier-Hein, L., M¨uller-Stich, B.P., Speidel, S., Distler, M., Weitz, J., Wagner, M.: A surgical activity model of laparoscopic cholecystectomy for co-operation with collaborative robots. Surgical Endoscopy 38(8), 4316–4328 (2024) https://doi.org/10.1007/s00464-024-10958-w. 
\item[13] Rueckert, T., Rauber, D., Maerkl, R., Klausmann, L., Yildiran, S.R., Gutbrod,
M., Nunes, D.W., Moreno, A.F., Luengo, I., Stoyanov, D., Toussaint, N., Cho, E., Kim, H.B., Choo, O.S., Kim, K.Y., Kim, S.T., Arantes, G., Song, K., Zhu, J., Xiong, J., Lin, T., Kikuchi, S., Matsuzaki, H., Kouno, A., Manesco, J.R.R., Papa, J.P., Choi, T.-M., Jeong, T.K., Park, J., Alabi, O., Wei, M., Vercauteren, T., Wu, R., Xu, M., Wang, A., Bai, L., Ren, H., Yamlahi, A., Hennighausen, J.,
Maier-Hein, L., Kondo, S., Kasai, S., Hirasawa, K., Yang, S., Wang, Y., Chen, H., Rodrıguez, S., Aparicio, N., Manrique, L., Lyons, J.C., Hosie, O., Ayobi, N., Arbel´aez, P., Li, Y., Al Khalil, Y., Nasirihaghighi, S., Speidel, S., Rueckert, D., Feussner, H., Wilhelm, D., Palm, C.: Comparative validation of surgical phase recognition, instrument keypoint estimation, and instrument instance segmentation in endoscopy: Results of the phakir 2024 challenge. Medical Image Analysis 109, 103945 (2026) https://doi.org/10.1016/j.media.2026.103945. 
\end{itemize}


\vspace{4mm}
We have also checked the entire reference list to ensure compliance with the journal guidelines.



\vspace{4mm}

\end{document}
